%=============================================================
% Paper 11: Acemoglu, Akcigit, Alp, Bloom & Kerr (2018), AER
%=============================================================
\chapter{Acemoglu, Akcigit, Alp, Bloom \& Kerr (2018)}

\begin{paperbox}{Innovation, Reallocation, and Growth}
\textbf{Authors:} Daron Acemoglu, Ufuk Akcigit, Harun Alp,
Nicholas Bloom, William Kerr\\[3pt]
\textbf{Journal:} \textit{American Economic Review} 108(11): 3450--3491\\[3pt]
\textbf{Year:} 2018 \quad\textbf{Field:} Macroeconomics; Endogenous growth;
Applied microeconomics\\[3pt]
\textbf{Classification:} Structural (endogenous growth model + SMM estimation)
\end{paperbox}

%------------------------------------------------------------
\section{Research Question}

What are the implications of different types of industrial policies---R\&D
subsidies, operating cost subsidies, and entry subsidies---for innovation,
reallocation, and growth in an economy where firms are heterogeneous in their
innovative capacity? Can \emph{taxing} the continued operation of incumbents
improve welfare by freeing skilled labour for R\&D by high-type firms?

%------------------------------------------------------------
\section{Motivation}

Industrial policies subsidising incumbent firms are pervasive worldwide
(EU: EUR~1.18 trillion in 2010, or 9.6\% of GDP). Yet existing models
either lack endogenous exit (and cannot model exit-targeting policies) or
lack firm heterogeneity in innovative capacity (and cannot model R\&D input
misallocation). The paper fills this gap: it shows the key market failure is
the misallocation of \emph{R\&D inputs} (skilled labour) between high- and
low-type firms, and demonstrates that operation taxes---not R\&D
subsidies---are the most effective instrument for positive selection.

%------------------------------------------------------------
\section{Main Contribution}

\begin{enumerate}
  \item \textbf{Model:} A quality-ladder endogenous growth model with two-type
    firm heterogeneity ($\theta_h > \theta_l$), endogenous exit via
    obsolescence, creative destruction, and type transitions
    (high$\to$low at rate $\nu$). Franchise values solved in closed form.

  \item \textbf{Structural estimation:} 8 parameters estimated by SMM
    targeting 18 moments from U.S.\ Census microdata on innovative
    manufacturing firms (LBD, CMF, NSF~R\&D, NBER Patents), 1987--1997.
    Covers 98\% of U.S.\ industrial R\&D.

  \item \textbf{Policy analysis (Table~12):} Incumbent R\&D subsidy = +0.63\%
    welfare; operation subsidy = $-$0.20\%; entry subsidy = $-$1.64\%;
    optimal operation \emph{tax} (69\%) = \textbf{+1.40\%}; social planner = +4.47\%.
\end{enumerate}

%------------------------------------------------------------
\section{Relationship to the Literature}

Extends Klette--Kortum (2004) and Lentz--Mortensen (2008) by adding:
(i)~endogenous exit via obsolescence; (ii)~two-type firm heterogeneity with
high-to-low transitions; (iii)~fixed costs of operation. Without (i) there is
no exit margin to tax; without (ii) all firms are homogeneous and R\&D
subsidies are not differentially distortionary. Also relates to the
misallocation literature (Hsieh--Klenow 2009; Restuccia--Rogerson 2008),
but focuses on misallocation of \emph{R\&D inputs} rather than production
inputs.

%------------------------------------------------------------
\section{Theoretical Framework and Economic Mechanism}

\subsection*{Preferences and Technology}

CRRA preferences over a CES consumption aggregate:
\begin{equation}
  U_0 = \int_0^\infty e^{-\rho t}\,\frac{C(t)^{1-\vartheta}-1}{1-\vartheta}\,dt,
  \qquad C(t) = \left(\int_{\mathcal{J}(t)} c_j(t)^{\frac{\varepsilon-1}{\varepsilon}}dj\right)^{\frac{\varepsilon}{\varepsilon-1}}.
  \label{eq:ace_prefs}
\end{equation}
Product~$j$ is produced by its leading-edge monopolist at cost $w_u/\hat{q}_j$
($\hat{q}_j \equiv q_j/w_u$), generating equilibrium profit:
\begin{equation}
  \pi(\hat{q}_j) = \Pi\,\hat{q}_j^{\varepsilon-1}, \qquad
  \Pi \equiv \frac{1}{\varepsilon-1}\left(\frac{\varepsilon-1}{\varepsilon}\right)^\varepsilon.
  \label{eq:ace_profit}
\end{equation}

\subsection*{Firm Type Dynamics}

Entrants draw $\theta=\theta_h$ with probability $\alpha$ and $\theta=\theta_l$
with probability $1-\alpha$. High-type firms transition to low-type at rate
$\nu>0$; low-type is absorbing. A type-$k$ firm with $n$ product lines and
$h$ skilled R\&D workers innovates at flow rate:
\begin{equation}
  X_f = \theta_k^\gamma\, n^\gamma\, h^{1-\gamma}, \qquad
  C(x,n,\theta) = w_s\,n\,G(x,\theta), \quad
  G(x,\theta) = x^{\frac{1}{1-\gamma}}\,\theta^{-\frac{\gamma}{1-\gamma}}.
  \label{eq:ace_innovation}
\end{equation}

\subsection*{Value Functions (Franchise Value Additivity)}

\textbf{Lemma 1:} The firm value is additive across product lines:
$\tilde{V}_k(\hat{\mathcal{Q}}) = \sum_{\hat{q}\in\hat{\mathcal{Q}}}\Upsilon_k(\hat{q})$,
where the franchise value $\Upsilon_k(\hat{q})$ satisfies the ODE
(Lemma~2):
\begin{equation}
  (r+\tau+\phi+(\varepsilon-1)g)\,\Upsilon_l(\hat{q})
  - \frac{\partial\Upsilon_l}{\partial\hat{q}}\frac{\partial\hat{q}}{\partial w_u}\dot{w}_u
  = \Pi\hat{q}^{\varepsilon-1} - \tilde{w}_s\phi + \Omega_l,
  \qquad \hat{q}>\hat{q}_{l,\min},
  \label{eq:ace_bellman_low}
\end{equation}
with closed-form solution (Proposition~1). The exit threshold is:
\begin{equation}
  \hat{q}_{k,\min} = \left(\frac{\tilde{w}_s\phi - \Omega_k}{\Pi}\right)^{\frac{1}{\varepsilon-1}},
  \quad k\in\{h,l\}.
  \label{eq:ace_exit}
\end{equation}

\subsection*{Growth and Key Market Failure}

Aggregate growth equals $g = \lambda\tau$ (Proposition~2; standard quality
ladder). The key market failure is the underinvestment in R\&D relative to the
social optimum: because future innovators build on current quality
($q_j(t+) = q_j + \lambda\bar{q}$), current innovation creates a positive
spillover. This underinvestment is especially costly from the perspective of
high-type firms, whose marginal product of R\&D is higher. The social planner
wishes to shift skilled labour from low-type firm operations to high-type R\&D.

\textbf{Why operation taxes work and R\&D subsidies fail:}
\begin{itemize}
  \item Low-type firms are disproportionately near the obsolescence threshold
    $\hat{q}_{l,\min}$ (Figure~2 in the paper).
  \item An operation tax raises $\hat{q}_{k,\min}$, forcing exit of low-type
    product lines near the margin. This frees skilled labour for R\&D by
    high-type firms, generating strong positive selection.
  \item R\&D subsidies raise demand for skilled labour from \emph{both} types,
    raising the skilled wage and crowding out entrant R\&D without improving
    the composition of active firms.
\end{itemize}

%------------------------------------------------------------
\section{Data}

\begin{itemize}
  \item \textbf{LBD:} Annual employment for all private-sector
    establishments, 1976--present.
  \item \textbf{CMF:} Quinquennial detailed plant/firm records; years
    1987, 1992, 1997.
  \item \textbf{NSF R\&D Survey (RAD):} Annual/biannual firm-level R\&D
    expenditures; certainty threshold $>\$1$m R\&D.
  \item \textbf{NBER Patent Database:} USPTO patents matched to LBD by
    name/location.
  \item \textbf{Sample:} ``Continuously innovative'' firms (R\&D or
    patenting in every CMF period). 17,055 obs.\ from 9,835 firms; 98\%
    of U.S.\ industrial R\&D; 50\% of manufacturing employment.
\end{itemize}

%------------------------------------------------------------
\section{Empirical Strategy}

SMM targeting 18 moments (firm exit rates by age/size, size transition
rates, employment and sales growth by age/size, R\&D intensity,
5-year entrant employment share, fixed cost/R\&D labour ratio, aggregate
growth). Aggregate growth is weighted $5\times$ the micro-moments.
Bootstrap standard errors (1,000 replications, stratified by firm age,
size, year, industry).

\subsection*{Computational Algorithm}

Solve the stationary equilibrium as a fixed point of six aggregate
variables:
\begin{equation}
  \{\tilde{w}_s,\,\Phi_h,\,\Phi_l,\,\hat{q}^{\,\text{avg}},\,
  \mathbb{E}[\Upsilon_h(\hat{q}+\lambda\hat{q}^{\,\text{avg}})],\,
  \mathbb{E}[\Upsilon_l(\hat{q}+\lambda\hat{q}^{\,\text{avg}})]\}.
  \label{eq:ace_fp}
\end{equation}
$2^{17}$ firms simulated for 10,000 periods (each period $= 0.02$ years)
to generate model-implied moments.

%------------------------------------------------------------
\section{Identification Strategy}

8 parameters are estimated from 18 moments (overidentified by 10). Each
parameter is primarily identified by specific moments:
\begin{itemize}
  \item $\phi$ (fixed cost): ratio of fixed-cost to R\&D workers.
  \item $\alpha$, $\nu$: age gradient in exit rates and R\&D intensity
    (most young firms are high-type; transitions to low-type raise exit
    rates for older cohorts).
  \item $\lambda$: aggregate growth rate.
  \item $\theta_h/\theta_l$: ratio of R\&D intensity between large-old
    and small-young firms.
  \item $\phi$ (exogenous exit): exit rates of large-old firms (unlikely
    to exit endogenously).
\end{itemize}

%------------------------------------------------------------
\section{Main Results}

\subsection*{Parameter Estimates (Table~2)}

\begin{center}
\small
\begin{tabular}{llc}
\toprule
\textbf{Parameter} & \textbf{Description} & \textbf{Value} \\
\midrule
$\phi$ & Fixed cost of operation & 0.216 \\
$\theta_h$ & High-type innovative capacity & 1.751 \\
$\theta_l$ & Low-type innovative capacity & 1.391 \\
$\theta_E$ & Entrant innovative capacity & 0.024 \\
$\alpha$ & P(high-type entrant) & 0.926 \\
$\nu$ & High$\to$low transition rate & 0.206/year \\
$\lambda$ & Innovation step size & 0.132 \\
$\phi$ (exog.\ exit) & Exogenous destruction rate & 0.037 \\
\bottomrule
\end{tabular}
\end{center}

Key implications: high-type firms are 26\% more innovative than low-type
($\theta_h/\theta_l\approx 1.26$); 93\% of entrants start as high-type; 21\%/year
probability of transitioning to low-type --- strong negative selection as firms age.

\subsection*{Baseline Economy (Table~4)}

Equilibrium growth: $g = 2.26\%$; creative destruction: $\tau = 17.2\%$;
R\&D employment share: $L_{R\&D}/L_S = 19.9\%$; product line shares:
$\Phi_h = 6.3\%$, $\Phi_l = 55.0\%$, $\Phi_{np} = 38.7\%$ (inactive).

\subsection*{Policy Experiments (Table~12)}

\begin{center}
\begin{tabular}{lcc}
\toprule
\textbf{Policy} & \textbf{Growth} & \textbf{Welfare} \\
\midrule
Baseline & 2.26\% & 100.00 \\
Incumbent R\&D subsidy (14\%, 1\% GDP) & 2.34\% & 100.63 \\
Operation subsidy (4\%, 1\% GDP) & 2.24\% & 99.80 \\
Entry subsidy (65\%, 1\% GDP) & 2.25\% & 98.36 \\
\textbf{Optimal operation tax ($-$69\%)} & \textbf{2.54\%} & \textbf{101.40} \\
Social planner & 2.94\% & 104.47 \\
\bottomrule
\end{tabular}
\end{center}

Social planner raises low-type exit threshold ($\hat{q}_{l,\min}$: $1.47\to 2.40$),
lowers high-type threshold ($\hat{q}_{h,\min}$: $1.30\to 0.28$), and raises
$\Phi_h/\Phi_l$ from 0.11 to 7.93.

%------------------------------------------------------------
\section{Economic Magnitude and Interpretation}

The welfare gain from the optimal operation tax (+1.40\%) is $2.2\times$ that
of the incumbent R\&D subsidy (+0.63\%), despite both costing 1\% of GDP.
The social planner achieves a 4.47\% welfare gain and raises growth by 0.68~pp
($2.26\to 2.94\%$), showing the equilibrium is substantially inefficient. The
operation tax is ``the most effective implementable policy'': by targeting firms
near the exit margin---disproportionately low-type---it achieves selection
effects that uniform R\&D subsidies cannot.

%------------------------------------------------------------
\section{Mechanisms}

The central mechanism operates through \textbf{reallocation of skilled labour}:
\begin{enumerate}
  \item Too many skilled workers cover fixed costs of low-type firm operations;
    too few perform R\&D.
  \item Operation tax raises $\hat{q}_{l,\min}$ $\Rightarrow$ forces exit of
    low-type product lines near the margin $\Rightarrow$ frees skilled labour.
  \item Freed skilled labour flows to R\&D by high-type firms (higher marginal
    product) $\Rightarrow$ positive selection: $\Phi_h/\Phi_l$ rises.
  \item R\&D subsidies fail because they increase skilled labour demand from
    \emph{both} types, raising the skilled wage and crowding out entrant R\&D.
\end{enumerate}

%------------------------------------------------------------
\section{Robustness Checks}

\begin{itemize}
  \item Results replicated with employment-weighted moments (vs.\ unweighted
    baseline); including non-innovative firms (dropping R\&D moments);
    excluding mergers and acquisitions.
  \item Model variants: fixed costs using skilled + unskilled labour; factor
    reallocation costs; more than two firm types; endogenous skill supply.
  \item Non-targeted moments: product line distribution (CMF Product
    Trailers), growth persistence by firm type, management scores from
    MOPS all closely matched.
\end{itemize}

%------------------------------------------------------------
\section{Limitations}

\begin{enumerate}
  \item Sample restricted to continuously innovative manufacturing firms
    (2\% of firms; 98\% of R\&D); conclusions may not generalise.
  \item No transitional dynamics between stationary equilibria; welfare
    gains exclude adjustment costs.
  \item Closed economy; no trade, multinationals, or international
    knowledge flows.
  \item Skilled labour relocates instantaneously; in reality R\&D labour
    is highly specialised and slow to reallocate.
  \item Two-type approximation; a richer type space would allow more
    nuanced policy targeting.
\end{enumerate}

%------------------------------------------------------------
\section{Policy Implications}

\begin{itemize}
  \item \textbf{R\&D subsidies are relatively ineffective.} +0.63\%
    welfare for 1\% GDP cost --- less than half the gain from an
    operation tax.
  \item \textbf{Taxing incumbent operations is powerful.} A 69\% tax on
    fixed costs (8\% of revenues on average) raises growth from 2.26\%
    to 2.54\% and welfare by 1.4\%.
  \item \textbf{Operation subsidies are harmful.} They increase low-type
    product line shares and depress R\&D; welfare falls.
  \item \textbf{Industrial policy should discriminate by exit margin, not
    by innovation rate.} Policies that target the exit margin achieve
    positive selection effects that R\&D subsidies cannot.
\end{itemize}

%------------------------------------------------------------
\section{Overall Assessment}

A technically demanding, ambitious paper that marries endogenous growth theory,
structural estimation, and policy analysis in a unified framework. The central
result---operation taxes dominate R\&D subsidies---is counterintuitive,
well-motivated, and robustly derived. The model is tractable (franchise values
in closed form), well-estimated (18 moments, strong non-targeted fit), and
policy analysis is comprehensive.

\textbf{Quality: Excellent. A major contribution to endogenous growth, firm
dynamics, and policy analysis.}

%------------------------------------------------------------
\section{Relevance to My Research}

Directly relevant to Canadian RDC research on industrial policy effectiveness.
The structural SMM approach is applicable to Canadian Innovation Survey data
linked to LEAP/T2 corporate tax records. The central question---whether SR\&ED
tax credits disproportionately benefit low-type vs.\ high-type Canadian
firms---would be a natural application with first-order policy implications.

%------------------------------------------------------------
\section{Potential Research Extensions}

\begin{itemize}
  \item Apply the Acemoglu et al.\ framework to Canadian data (LEAP/T2/
    Innovation Survey); test whether SR\&ED credits disproportionately
    benefit low-type firms as predicted.
  \item Extend to a small open economy to capture import competition as a
    mechanism for forcing low-type firm exit (analogous to the operation tax).
\end{itemize}

%------------------------------------------------------------
\section{Research Gaps}

\begin{itemize}
  \item The political economy of why operation taxes are not observed in
    practice, despite their welfare superiority, is unmodelled.
  \item No analysis of heterogeneity in policy effects across industries
    with different innovation intensities (high-tech vs.\ low-tech).
\end{itemize}

%------------------------------------------------------------
\section{Methodological Lessons}

\begin{enumerate}
  \item The franchise value additivity result (Lemma~1) reduces the state
    space from the full product portfolio to a scalar per product line,
    enabling closed-form solutions even in a complex GE environment.
  \item SMM with 18 moments and 8 parameters is overidentified; matching
    non-targeted moments provides external validation that goes beyond the
    targeted moments.
  \item The policy comparison framework---all interventions costing exactly
    1\% of GDP---enables clean welfare comparison across fundamentally
    different policy instruments.
\end{enumerate}

%------------------------------------------------------------
\section*{Five-Sentence Summary}

Acemoglu et al.\ (2018) develop a quality-ladder endogenous growth model
featuring high- and low-type firms (differing in innovative capacity) with
endogenous exit, creative destruction, and type transitions, estimated by SMM
on 18 moments from U.S.\ Census Bureau microdata covering 98\% of U.S.\
industrial R\&D. The model's key market failure is the misallocation of skilled
labour between low-type firm operations (generating little social value) and
R\&D by high-type firms (generating large positive spillovers), leading to
underinvestment in socially valuable innovation. Policy analysis shows that
incumbent R\&D subsidies are relatively ineffective (+0.63\% welfare for 1\%
of GDP) because they benefit both firm types indiscriminately, while a 69\%
operation tax generates a strong positive selection effect (+1.40\% welfare)
by disproportionately inducing exit among low-type firms and freeing skilled
labour for high-type R\&D. The social planner, controlling type-specific exit
thresholds and R\&D rates, achieves +4.47\% welfare and raises growth from
2.26\% to 2.94\% by dramatically raising the low-type exit threshold and
redirecting skilled labour to high-type R\&D. The central policy lesson is that
the standard tool of industrial policy---R\&D subsidies---is not merely
ineffective but is dominated by policies that accelerate exit of
low-productivity incumbents.

\vspace{6pt}
\begin{quickref}
\begin{tabularx}{\linewidth}{@{}lX@{}}
\textbf{Key contribution} &
  Incumbent operation taxes dominate R\&D subsidies for welfare: taxes fall
  disproportionately on low-type firms near the exit margin, freeing skilled
  labour for R\&D by high-type firms with a 1.4\% welfare gain vs.\ 0.6\%
  from R\&D subsidies. \\[4pt]
\textbf{Key weakness} &
  Sample restricted to continuously innovative manufacturing firms (2\% of
  all firms); no transitional dynamics; closed economy; instantaneous skilled
  labour reallocation. \\[4pt]
\textbf{Most interesting gap} &
  The political economy of why operation taxes are not observed in practice
  despite their welfare superiority---incumbents lobby against exit-inducing
  policies even when they reduce aggregate growth. \\[4pt]
\textbf{Publishable extension} &
  Apply the framework to Canadian data (LEAP/T2/Innovation Survey); test
  whether SR\&ED credits disproportionately benefit low-type vs.\ high-type
  Canadian firms, and estimate the implied welfare cost of the current design.
\end{tabularx}
\end{quickref}
